\section*{第 \ref{ch:b3:developmental-genetics} 章——发育遗传学与形态演化}

\begin{solution}{exo:b3:developmental-genetics:1}
母体效应：\emph{bicoid}，突变型母体所产的胚胎缺少头部
和胸部。\omterm{def:b3:developmental-genetics:hierarchy}{裂隙基因}：\emph{Krüppel}，连续一段体节
缺失。\omterm{def:b3:developmental-genetics:hierarchy}{成对规则基因}：\emph{fushi tarazu}（或 \emph{even-skipped}），每
隔一个体节缺失。\omterm{def:b3:developmental-genetics:hierarchy}{体节极性基因}：\emph{hedgehog}（或
\emph{engrailed}），每个体节的一部分被替换为
其余部分的镜像。
\end{solution}

\begin{solution}{exo:b3:developmental-genetics:2}
一种基因，其产物由母体沉积于卵中，
使胚胎的表型随母体基因型而定。其胚胎缺少头部和
胸部，两端带有后部结构，胚胎死亡。
父本的野生型等位基因存在于受精卵中，但\emph{bicoid} mRNA
是在卵发生期间由母体滋养细胞合成，并在
受精前锚定于卵的一极；受精卵自身的拷贝要在数小时后才被转录，
那时梯度形成的时间早已过去。
\end{solution}

\begin{solution}{exo:b3:developmental-genetics:3}
\omterm{def:b3:developmental-genetics:hox}{共线性}：Hox 基因在染色体上的排列顺序，与它们
在身体上的表达域顺序一致。后部优先：
若多个基因同时表达，最靠后的那个决定身份。
缺失 \emph{Hoxc8} 的小鼠表现出前部转化——其
第一节腰椎获得了胸椎的身份，并长出一根
肋骨。头部的 \emph{Antennapedia} 使触角转变为腿，
即第二胸节的身份。
\end{solution}

\begin{solution}{exo:b3:developmental-genetics:4}
移植的背唇\omterm{def:b3:developmental-genetics:organiser}{诱导}出第二个完整的轴，该轴主要由宿主
细胞构成：移植体指导了其邻近细胞——即\omterm{def:b3:developmental-genetics:organiser}{组织者}。它的
分子是分泌性 BMP 拮抗剂（Chordin、Noggin、Follistatin）。
外胚层在无 BMP 信号时成为神经组织，在有 BMP 时成为表皮，
因此神经组织是外胚层在未接收任何指令时的默认状态：
\omterm{def:b3:developmental-genetics:organiser}{组织者}通过沉默一条指令来发挥作用。
\end{solution}

\begin{solution}{exo:b3:developmental-genetics:5}
$k = \ln 2/(35\times 60) = \qty{3.3e-4}{s^{-1}}$; $\lambda = \sqrt{4/
3.3\times 10^{-4}} = \qty{110}{\micro m}$；比值 $\mathrm{e}^{250/110}
= \mathrm{e}^{2.27} = 9.7$.
\end{solution}

\begin{solution}{exo:b3:developmental-genetics:6}
四个拷贝使 $c_{0}$ 加倍：边界向后移动 $\lambda\ln 2 =
\qty{69}{\micro m}$ 至 \qty{309}{\micro m}；一个拷贝使其减半：边界向前
移动 \qty{69}{\micro m} 至 \qty{171}{\micro m}。每次位移
占卵长的 \qty{14}{\%}——比观察值大，这本身就是
证据，说明存在能够减弱剂量依赖性的机制。
\end{solution}

\begin{solution}{exo:b3:developmental-genetics:7}
$\delta x = \lambda\,\delta c/c$：对于 \qty{8}{\micro m}，$\delta c/c =
8/100 = \qty{8}{\%}$。泊松计数 $N$ 的相对误差为
$1/\sqrt{N}$，故 $N = 1/0.08^{2} \approx 160$ 个分子须被计数，
——约占细胞核内分子数的百分之几，或一次
随时间累积的计数。
\end{solution}

\begin{solution}{exo:b3:developmental-genetics:8}
$k = \ln 2/35 = \qty{0.0198}{min^{-1}}$：$1 - \mathrm{e}^{-kt}$ 给出
$0.70$ 在 \qty{60}{min}，$0.83$ 在 \qty{90}{min}，$0.91$ 在
\qty{120}{min}，$0.95$ 在 \qty{150}{min}。梯度在被读取期间仍在
以百分之几的幅度上升，这会使边界位移
$\lambda\ln(1/0.9) \approx \qty{10}{\micro m}$，约相当于一个细胞核的宽度，在
读取时间窗口内。胚胎必须要么在固定时间读取，要么使用
对整体水平不敏感的读出机制。
\end{solution}

\begin{solution}{exo:b3:developmental-genetics:9}
第二个 ZPA 产生镜像重复，4-3-2-2-3-4：
前部细胞现在感受到高浓度的 Shh，并发育为后部指。一个
释放少量 Shh 的珠子会产生部分重复——多出一个前部
指 2，例如 2-2-3-4——因为前部细胞只感受到较低浓度，
该浓度恰好决定指 2 的身份。
\end{solution}

\begin{solution}{exo:b3:developmental-genetics:10}
特解为 $s/k$；齐次解为 $\cosh(x/\lambda)$，
$\sinh(x/\lambda)$；$0$ 处通量为零，故舍去 $\sinh$；$c(L) = 0$ 确定
振幅：$c(x) = (s/k)\bigl[1 - \cosh(x/\lambda)/\cosh(L/\lambda)\bigr]$。
该分布在 $x = 0$ 附近为平台，在汇处宽度为 $\lambda$ 的边界层内
降至零——它不是来自一点的指数分布，
而是带有边缘的平坦场；浓度最高处
离汇最远。对于 $L \ll \lambda$，$\cosh u \approx 1 +
u^{2}/2$，且 $c \approx s(L^{2} - x^{2})/(2D)$：这是与 $k$ 无关的抛物线，
因为分子在衰变之前就已到达汇。
\end{solution}

\begin{solution}{exo:b3:developmental-genetics:11}
$x_{1} = \lambda\ln(1/0.3) = 1.20\lambda$ 且 $x_{2} = \lambda\ln(1/0.08)
= 2.53\lambda$：第一区域宽 $1.20\lambda$，第二区域
宽 $1.32\lambda$，二者均与 $c_{0}$ 无关，整体随之平移。
第三区域 $L - 2.53\lambda$ 吸收了卵长的全部变化。
例如，在卵长 \qty{450}{\micro m} 且 $\lambda = \qty{100}{\micro m}$ 的卵中，
腹部为 \qty{197}{\micro m}，在卵长 \qty{550}{\micro m} 的卵中
为 \qty{297}{\micro m}——图式并不随比例缩放。可能的机制：第二个
梯度来自后极（Nanos、Caudal 或末端
系统），与 Bicoid 共同读取，使位置由
二者的比值决定；或者在\omterm{thm:b3:developmental-genetics:gradient}{形态发生素}浓度低的地方产生一种“扩张子”分子，
它使 $\lambda$ 增大，直至梯度充满整个卵。
\end{solution}

\begin{solution}{exo:b3:developmental-genetics:12}
编码序列的改变会使蛋白质在其作用的任何地方都发生改变：\emph{Pitx1} 还
构建\omterm{def:b3:endocrinology:axes}{垂体}和颌部，\emph{sonic hedgehog} 也参与神经管、肠道和牙齿的模式形成，
因此蛋白质一旦受损，要么致死，要么
产生多效性致残。增强子驱动表达于某一组织，
且仅在某一时间；删除它则去除该组织本应形成的结构，
而蛋白质的其他一切作用保持完好。增强子是
模块化的，且常有部分冗余，因此中间步骤是可行的，
而同一路径——开关丢失——已在
多个刺鱼湖泊中独立重现。
\end{solution}

\begin{solution}{pb:b3:developmental-genetics:1}
\textbf{1.} $k = \ln 2/(35\times 60) = \qty{3.3e-4}{s^{-1}}$; $1/k =
\qty{3030}{s} = \qty{50}{min}$.
\textbf{2.} $\lambda = \sqrt{4/3.3\times 10^{-4}} = \qty{110}{\micro m}$;
卵长为 $4.5\lambda$。
\textbf{3.} $\mathrm{e}^{4.5} \approx 90$，从一极到另一极；
$\mathrm{e}^{8.5/110} = 1.08$，每个细胞核处变化 \qty{8}{\%}。
\textbf{4.} $J = 50\times\sqrt{4\times 3.3\times 10^{-4}} = 50\times
0.036 = \qty{1.8}{nM\,\micro m\,s^{-1}}$。1 nM 约为 $0.6$ 个分子
每 \unit{\micro m^3}，故每 \unit{\micro m^2} 约 $1.1$ 个分子
每秒。
\textbf{5.} 细胞核体积 $\tfrac{4}{3}\pi\times 27 = \qty{113}{\micro
m^3}$。当 \qty{50}{nM} 时，每 \unit{\micro m^3} 含 $30$ 个分子：
\num{3400} 个分子。在卵中部，$c = 50\,\mathrm{e}^{-2.27} =
\qty{5.2}{nM}$，约 \num{350} 个分子。
\textbf{6.} $kt = 1.19$（\qty{60}{min}）：$0.70$；$kt = 2.38$
（\qty{120}{min}）：$0.91$。梯度与稳态的偏差在 \qty{10}{\%} 以内
读取时——仍在上升，增量足以使边界移动
$\lambda\ln(1/0.91) \approx \qty{10}{\micro m}$，即一个细胞核。
\textbf{7.} $x_{1} = 110\ln(1/0.3) = \qty{132}{\micro m}$
(\qty{26}{\%}); $x_{2} = 110\ln(1/0.08) = \qty{278}{\micro m}$
(\qty{56}{\%}).
\textbf{8.} 四个拷贝：$c_{0} = \qty{100}{nM}$；两个边界均
向后移动 $\lambda\ln 2 = \qty{76}{\micro m}$，分别到达 $208$ 和
\qty{354}{\micro m}。六个拷贝：$\lambda\ln 3 = \qty{121}{\micro m}$，
分别到达 $253$ 和 \qty{399}{\micro m}。
\textbf{9.} 一个拷贝：向前移动 \qty{76}{\micro m}，分别到达 $56$ 和
\qty{202}{\micro m}。头部区域由 $132$ 缩小至
\qty{56}{\micro m}；胸部保持宽度（\qty{146}{\micro m}）并
向前移动；腹部由 $222$ 增长至 \qty{298}{\micro m}。
\textbf{10.} $\delta x = 110\times 0.1 = \qty{11}{\micro m}$, $1.3$
个细胞核间距。
\textbf{11.} 单个细胞核：$\delta c/c = 8.5/110 = \qty{7.7}{\%}$。要达到
\qty{10}{\%}，需 $n = (10/7.7)^{2} \approx 1.7$：两次独立读取
（或取相邻两个细胞核的平均值）即已足够。
\textbf{12.} $1/\sqrt{350} = \qty{5.3}{\%}$：单次瞬时计数
已比假设的 \qty{10}{\%} 更精确，因此
观测到的噪声并非仅来自计数噪声，而是来自结合
动力学与读出机制——而对时间求平均
则能改善这一点。
\textbf{13.} $278/450 = \qty{62}{\%}$ 与 $278/550 = \qty{51}{\%}$，占
卵长的比例：相差 \qty{11}{\%}，相当于六个细胞核。该模式不
随固定的 $\lambda$ 而缩放。
\textbf{14.} 若 $D \propto L^{2}$、$k$ 固定，则 $\lambda \propto L$ 且
$x_{T}/L = (\lambda/L)\ln(c_{0}/T)$ 在每个卵中均相同，即完美
缩放。然而 $D$ 是分子与细胞质的性质，而非
卵的大小，因此按所述条件这一假设不太可能成立；真实胚胎中的缩放
只是部分的，并来自其他读出机制。
\textbf{15.} $c^{5}/(c^{5} + T^{5}) = 0.1$ 时 $c/T = 9^{-1/5} = 0.64$，而
$0.9$ 时 $c/T = 9^{1/5} = 1.55$：浓度相差 $2.4$ 倍，
$\Delta x = 110\ln 2.4 = \qty{97}{\micro m}$——相当于十一个细胞核，远
宽于观测到的两到三个细胞核宽的边界；梯度的
读出机制本身不够锐利。
\textbf{16.} 陡度使给定浓度更果断地映射为开或关，
过渡区因此更窄；但 \qty{10}{\%} 的
浓度误差仍会使阈值被越过的位置移动
$\lambda\,\delta c/c$：响应可以是完美的阶跃，
但阶跃仍落在错误的位置。
\textbf{17.} 若两者同时开启，则各自都会抑制对方，因此
至少有一个会被关闭；在协同抑制下，唯一
稳定状态是一开一关，而两者之间的边界是
开关，而非斜坡。
\textbf{18.} $\mathrm{d}\ln\lambda = \tfrac{1}{2}(0.02 - 0.06) = -0.02$
每度：$\lambda$ 下降 \qty{2}{\%}，约 \qty{10}{\%}，对应
\qty{5}{\celsius}，降至 \qty{99}{\micro m}；\emph{hunchback}
边界由 $278$ 移至 \qty{250}{\micro m}，向前移动三个细胞核的距离。
果蝇在 \qty{18}{} 与 \qty{29}{\celsius} 之间发育，比例正常，
因此必有补偿作用——某个源或读出机制，
其随温度的变化方向与之相反。
\textbf{19.} $2^{8} = 256$ 种编码，用到十四种。在后部优势下
只有最靠后的开启基因有作用，因此至多九种
不同身份（八个基因之一，或无）：该编码是排序，而非
二进制字。
\textbf{20.} 第三胸节变成了第二胸节：\emph{Ubx}
抑制了胸节默认运行的翅发育程序，
去除抑制因子便释放了祖先状态——四翅，
正如果蝇的四翅祖先。
\textbf{21.} 鹅的颈区延伸至第 22 节椎骨：一条长
颈，由多节椎骨构成，小鼠则为七节。两者基因相同；
其表达的前界向后移动——一种调控上的
变化，即 Hox 基因在何处被开启。
\textbf{22.} 现在每个边缘都是一个源；Shh 在两端最高
（第 4 指），向中心下降（先为第 3 指，再为第 2 指），且从未达到
决定第 1 指的低水平，因此中部有两个第 2 指
背靠背，排列为 4-3-2-2-3-4。
\textbf{23.} \emph{Pax6} 是一种保守的主调控因子，能
启动宿主所携带的任何眼发育程序：果蝇构建的是果蝇眼，
因此小鼠的基因并不携带小鼠眼的信息，只携带
“在此处构建眼”的指令。它显示的是基因网络的\omterm{def:b3:developmental-genetics:evodevo}{深层同源}，
而非器官的同源：复眼与照相机眼
各自独立地由一个具有 \emph{Pax6}
光感受器斑块的祖先演化而来。
\textbf{24.} 五指横跨 \qty{2}{mm}，对应的波长为
\qty{0.4}{mm}；六指则需 \qty{0.33}{mm}，缩短了六分之一。
由于 $\lambda \propto \sqrt{D_{\mathrm{inh}}/k_{\mathrm{inh}}}$，提高
抑制物的衰减速率 \qty{44}{\%}，或使其扩散降低
\qty{31}{\%}；在手板中，降低后部 \omterm{def:b3:developmental-genetics:hox}{Hox
基因}的剂量即可做到，而突变小鼠具有六、七或更多纤细的
指。
\textbf{25.} $\lambda \approx \qty{110}{\micro m}$；源的加倍
使每个边界移动 \qty{76}{\micro m}；浓度中的 \qty{10}{\%} 噪声
相当于 \qty{11}{\micro m}，约一个细胞核，
即位置误差。
\end{solution}
